% !PRINTED 192
\weberpage{192}
%<<CONT
so bestimmen, dass sowohl die Bedingung (13) als (14) befriedigt
wird. % [sic/wording] the common factor has two choices, satisfying (13) or (14), respectively
Demnach bekommen wir aus jeder Substitution $A'''$ zwei
orthogonale Substitutionen $A$, von denen die eine eigentlich, die
andere uneigentlich ist.

Die drei Substitutionen
\weq{19}{A,\ A',\ A''}
entsprechen sich also hiernach gegenseitig eindeutig, die Substitutionen

\weq{20}{A,\ A',\ A'''}
aber nur dann, wenn wir noch die Forderung stellen, dass $A$
eine eigentliche orthogonale Substitution sein soll.

Ist nun
\weq{21}{B,\ B',\ B''}
ein zweites System von Substitutionen der Form (19), so sind
auch $AB$, $A'B'$ zwei einander entsprechende Substitutionen, wie
unmittelbar aus der Bedeutung der $A'$ als transformirte Substitution
 der $A$ hervorgeht. Es ist aber noch nachzuweisen,
dass auch
\weq{22}{AB,\ A'B',\ A''B''}
ein zusammengehöriges System von der Form (19) ist, dass also
die Gruppen der $A, A', A''$ isomorph sind. Daraus ergiebt sich
dann von selbst aus der Beziehung, in der $A''$ und $A'''$ zu einander
 stehen, dass (bei Beschränkung auf eigentlich orthogonale
Substitutionen $A$) auch die Gruppe der $A'''$ damit isomorph ist.

Setzen wir, um dies zu beweisen:
\weq{23}{(\eta_1,\eta_2) = A''(\xi_1,\xi_2),\quad (\xi_1,\xi_2) = B''(\zeta_1,\zeta_2),}
so folgt
\weq{24}{(\eta_1,\eta_2) = A''B''(\zeta_1,\zeta_2).}

Aus $A''$, $B''$, $A''B''$ leiten wir nun nach (12) drei ternäre
Substitutionen $A'$, $B'$, $C'$ her. So erhalten wir nach (8):
\weq{25}{\begin{aligned}
(\sqrt{2}\,\eta_1\eta_2,\ \eta_1^2,\ \eta_2^2) &= A'(\sqrt{2}\,\xi_1\xi_2,\ \xi_1^2,\ \xi_2^2)\\
(\sqrt{2}\,\xi_1\xi_2,\ \xi_1^2,\ \xi_2^2) &= B'(\sqrt{2}\,\zeta_1\zeta_2,\ \zeta_1^2,\ \zeta_2^2)\\
(\sqrt{2}\,\eta_1\eta_2,\ \eta_1^2,\ \eta_2^2) &= C'(\sqrt{2}\,\zeta_1\zeta_2,\ \zeta_1^2,\ \zeta_2^2).
\end{aligned}}

Diese Formeln sind in Bezug auf $\zeta_1, \zeta_2$ identisch, und sie
müssen also richtig bleiben, wenn $\sqrt{2}\,\zeta_1\zeta_2, \zeta_1^2, \zeta_2^2$ durch drei
%CONT>>
